% problem-type: nhiet-hinh-chu-nhat
% order: 67
% Nguon: De thi MAT3365 (PTDHR Toan tin), cuoi ky, K68, Cau 2. (co dap an chinh thuc)

\section{Nhiệt trong hình chữ nhật --- năng lượng 2D \& thác triển/Poisson}

\subsection*{Đề bài ví dụ}

Xét bài toán biên --- ban đầu cho phương trình truyền nhiệt trong hình chữ nhật
$R=(-1,1)\times(0,\pi)$:
\[
  u_t = u_{xx} + u_{yy} + F(x,y,t), \qquad (x,y)\in R,\ t>0,
\]
với điều kiện biên
\[
  u_x(-1,y,t)=u_x(1,y,t)=0, \qquad u(x,0,t)=u(x,\pi,t)=0,
\]
và điều kiện ban đầu $u(x,y,0)=g(x,y)$.

\begin{enumerate}
\item[(a)] Dùng tích phân năng lượng $I(t)=\iint_R u^2\,dx\,dy$ chứng minh bài toán có duy nhất nghiệm.
\item[(b)] Với $F=0$, $g(x,y)=\sin^2(\pi x)\sin(y)$, bằng cách thác triển lên toàn mặt phẳng rồi dùng công thức Poisson, hãy giải bài toán.
\end{enumerate}

\emph{\small(Nguồn: Đề thi MAT3365 --- PTĐHR Toán tin, cuối kỳ, K68, Câu 2. Có đáp án chính thức.)}

\subsection*{Nhận ra dạng này khi nào?}

\begin{itemize}
\item Phương trình nhiệt trên miền \textbf{chữ nhật 2D} $R=(a,b)\times(c,d)$: $u_t=\Delta u+F$.
\item Biên cho trên \textbf{bốn cạnh}, mỗi cặp cạnh đối diện theo một biến: kiểu Dirichlet ($u=0$) hay Neumann ($u_x=0$/$u_y=0$).
\item Thường hỏi: (a) duy nhất nghiệm bằng năng lượng \emph{2D}; (b) giải tường minh.
\end{itemize}

\begin{meobox}
Mỗi biến tách riêng. Loại biên theo từng biến quyết định kiểu thác triển:
\begin{itemize}
\item Cặp cạnh \textbf{Dirichlet} ($u=0$) $\Rightarrow$ thác triển \textbf{lẻ}, tuần hoàn.
\item Cặp cạnh \textbf{Neumann} ($u_x=0$) $\Rightarrow$ thác triển \textbf{chẵn}, tuần hoàn.
\end{itemize}
Sau khi thác triển $\Rightarrow$ bài Cauchy trên toàn mặt phẳng, nhân kernel tách thành tích hai tích phân 1D.
\end{meobox}

\subsection*{Công thức \& phương pháp}

\begin{dieukienbox}
Dùng khi: nhiệt $u_t=\Delta u+F$ trên hình chữ nhật, biên Dirichlet/Neumann theo từng cặp cạnh; cần chứng minh duy nhất và/hoặc giải tường minh.
\end{dieukienbox}

\begin{nhobox}
\textbf{(1) Năng lượng 2D $\Rightarrow$ duy nhất.} Hiệu hai nghiệm $w$ có $F=0$, $w(\cdot,0)=0$. Đặt $I(t)=\iint_R w^2\ge0$:
\[
  I'(t)=2\!\iint_R w\,\Delta w
  =2\oint_{\partial R} w\,\partial_n w\,ds-2\!\iint_R|\nabla w|^2 .
\]
Biên: Dirichlet $\Rightarrow w=0$; Neumann $\Rightarrow \partial_n w=0$ $\Rightarrow$ số hạng biên triệt tiêu. Vậy $I'\le0$, $I(0)=0\Rightarrow I\equiv0\Rightarrow w\equiv0$.

\textbf{(2) Thác triển $\Rightarrow$ Cauchy toàn mặt phẳng.} Thác triển $u$ (và $g$) theo đúng đối xứng của biên. Nghiệm Cauchy nhiệt 2D tách biến:
\[
  u(x,y,t)=\Big[\tfrac{1}{\sqrt{4\pi t}}\!\int_{\mathbb R}\!\bar g_x(X)e^{-\frac{(x-X)^2}{4t}}dX\Big]
  \Big[\tfrac{1}{\sqrt{4\pi t}}\!\int_{\mathbb R}\!\bar g_y(Y)e^{-\frac{(y-Y)^2}{4t}}dY\Big].
\]

\textbf{(3) Làm trơn nhiệt của $\cos/\sin$:} dùng tích phân cos-Gauss (Phụ lục B),
\[
  \tfrac{1}{\sqrt{4\pi t}}\!\int_{\mathbb R}\!\cos(\beta X)e^{-\frac{(x-X)^2}{4t}}dX=e^{-\beta^2 t}\cos(\beta x),
\]
tương tự cho $\sin$. Kernel làm tắt mỗi mode tần số $\beta$ theo $e^{-\beta^2 t}$.
\end{nhobox}

\subsection*{Đề bài \& bài giải}

% ===============================================================
\subsubsection*{(a) Duy nhất nghiệm bằng năng lượng 2D}
% ===============================================================

\paragraph{Bước 1 --- lập bài toán cho hiệu hai nghiệm.}\leavevmode

\emph{Ý tưởng:} muốn chứng minh \textbf{duy nhất nghiệm}, giả sử có hai nghiệm $u_1,u_2$ rồi chứng minh hiệu $w:=u_1-u_2\equiv0$. Vì PDE và mọi điều kiện (biên, đầu) đều \textbf{tuyến tính} và \textbf{cùng dữ kiện}, hiệu khử sạch nguồn $F$ và dữ kiện ban đầu:
\[
  w_t=w_{xx}+w_{yy}\ \text{trên}\ R,\quad
  w_x(\pm1,y,t)=0,\quad w(x,0,t)=w(x,\pi,t)=0,\quad w(x,y,0)=0.
\]

\paragraph{Bước 2 --- định nghĩa năng lượng và tính đạo hàm.}\leavevmode

\textbf{Tích phân năng lượng} là đại lượng vô hướng $I(t)=\iint_R w^2\,dx\,dy\ge0$ đo ``độ lớn'' tổng cộng của $w$ tại thời điểm $t$. Hai tính chất chìa khoá: $I\ge0$ luôn, và $I(0)=0$ do $w(\cdot,0)=0$. Lấy đạo hàm dưới dấu tích phân (miền $R$ cố định) rồi dùng $w_t=\Delta w$:
\[
  I'(t)=\iint_R 2w\,w_t\,dx\,dy=2\iint_R w\,(w_{xx}+w_{yy})\,dx\,dy.
\]

\paragraph{Bước 3 --- tích phân từng phần (Green) để tách phần biên.}\leavevmode

Đẳng thức Green chuẩn (tích phân từng phần 2D) cho:
\[
  \iint_R w\,\Delta w\,dx\,dy=\oint_{\partial R} w\,\partial_n w\,ds-\iint_R |\nabla w|^2\,dx\,dy,
\]
với $\partial_n$ là đạo hàm pháp tuyến ngoài và $|\nabla w|^2=w_x^2+w_y^2$. Vậy
\[
  I'(t)=2\oint_{\partial R} w\,\partial_n w\,ds\;-\;2\iint_R(w_x^2+w_y^2)\,dx\,dy.
\]

\paragraph{Bước 4 --- triệt tích phân biên (xét từng cặp cạnh).}\leavevmode

Biên $\partial R$ gồm 4 cạnh; trên mỗi cạnh tích $w\,\partial_n w$ tự triệt:
\begin{itemize}
\item Cạnh $x=\pm1$ (Neumann): $\partial_n w=\pm w_x=0\Rightarrow w\,\partial_n w=0$.
\item Cạnh $y=0,\pi$ (Dirichlet): $w=0\Rightarrow w\,\partial_n w=0$.
\end{itemize}
Do đó số hạng biên triệt tiêu hoàn toàn:
\[
  I'(t)=-2\iint_R(w_x^2+w_y^2)\,dx\,dy\le0.
\]

\paragraph{Bước 5 --- kết luận từ $I\ge0$, $I(0)=0$, $I'\le0$.}\leavevmode

Ba điều này ép $0\le I(t)\le I(0)=0$ với mọi $t\ge0$, nên $I\equiv0$. Vì $w^2\ge0$ liên tục, tích phân $\iint w^2=0$ kéo theo $w\equiv0$ trên $R$. Vậy $u_1\equiv u_2$: \textbf{bài toán có duy nhất nghiệm}. \hfill$\square$

% ===============================================================
\subsubsection*{(b) Giải tường minh với $F=0$, $g=\sin^2(\pi x)\sin y$}
% ===============================================================

\emph{Chiến lược:} thác triển $u$ và $g$ từ hình chữ nhật $R$ lên toàn mặt phẳng $\mathbb R^2$ sao cho thoả \emph{tự động} cả 4 điều kiện biên; bài toán biên-đầu trở thành \textbf{bài Cauchy} trên $\mathbb R^2$ giải bằng công thức Poisson. Sau đó tách kernel Gauss thành tích hai tích phân 1D rồi tính từng cái.

\paragraph{Bước 1 --- thác triển: chọn đối xứng khớp từng loại biên.}\leavevmode

Ta cần \emph{chuyển} biên thành một \emph{tính chất đối xứng} của hàm thác triển $\bar u$, để biên tự thoả mà không phải áp đặt:
\begin{itemize}
\item Theo $y$ (biên Dirichlet $u=0$ tại $y=0,\pi$): thác triển \textbf{lẻ} qua mỗi đầu, rồi tuần hoàn chu kỳ $2\pi$. Một hàm lẻ tại $y=0$ tự bằng $0$, ổn với Dirichlet.
\item Theo $x$ (biên Neumann $u_x=0$ tại $x=\pm1$): thác triển \textbf{chẵn} qua mỗi đầu, rồi tuần hoàn chu kỳ $4$. Một hàm chẵn quanh $x=1$ thì $u_x(1,\cdot)=0$ tự động, ổn với Neumann.
\end{itemize}
(Quy tắc tổng quát: Dirichlet$\to$lẻ; Neumann$\to$chẵn --- xem Phụ lục C.)

May mắn $g(x,y)=\sin^2(\pi x)\sin(y)$ \emph{đã sẵn} đúng cả hai đối xứng:
\begin{itemize}
\item $\sin y$ lẻ và $2\pi$-tuần hoàn theo $y$ $\Rightarrow$ khớp Dirichlet ở $y=0,\pi$;
\item $\sin^2(\pi x)=\tfrac12(1-\cos(2\pi x))$ chẵn quanh mỗi $x=k$ ($k\in\mathbb Z$) và $1$-tuần hoàn $\Rightarrow$ khớp Neumann ở $x=\pm1$.
\end{itemize}
Do đó $\bar g(x,y)=\sin^2(\pi x)\sin(y)$ trên toàn $\mathbb R^2$.

\paragraph{Bước 2 --- công thức Poisson 2D \& kernel tách biến.}\leavevmode

Bài Cauchy $\bar u_t=\Delta\bar u$ trên $\mathbb R^2$, $\bar u(\cdot,0)=\bar g$ có nghiệm
\[
  \bar u(x,y,t)=\frac{1}{4\pi t}\!\iint_{\mathbb R^2}\bar g(X,Y)\,e^{-\frac{(x-X)^2+(y-Y)^2}{4t}}dX\,dY.
\]
Vì $\bar g(X,Y)=\sin^2(\pi X)\cdot\sin(Y)$ tách thành tích, và kernel Gauss cũng tách:
\[
  e^{-\frac{(x-X)^2+(y-Y)^2}{4t}}=e^{-\frac{(x-X)^2}{4t}}\,e^{-\frac{(y-Y)^2}{4t}},\qquad
  \frac{1}{4\pi t}=\frac{1}{\sqrt{4\pi t}}\cdot\frac{1}{\sqrt{4\pi t}},
\]
nên tích phân đôi rã thành tích hai tích phân 1D:
\[
  \bar u(x,y,t)=\underbrace{\Big[\tfrac{1}{\sqrt{4\pi t}}\!\int_{\mathbb R}\!\sin^2(\pi X)\,e^{-\frac{(x-X)^2}{4t}}dX\Big]}_{=:\,A(x,t)}
  \cdot\underbrace{\Big[\tfrac{1}{\sqrt{4\pi t}}\!\int_{\mathbb R}\!\sin(Y)\,e^{-\frac{(y-Y)^2}{4t}}dY\Big]}_{=:\,B(y,t)}.
\]

\paragraph{Bước 3 --- bổ đề: làm trơn nhiệt của $\cos/\sin$.}\leavevmode

Phụ lục B cho công thức (tích phân cos-Gauss): với mọi $\beta\in\mathbb R$,
\[
  \tfrac{1}{\sqrt{4\pi t}}\!\int_{\mathbb R}\!\cos(\beta X)\,e^{-\frac{(x-X)^2}{4t}}dX=e^{-\beta^2 t}\cos(\beta x),
\]
và tương tự thay $\cos\to\sin$. Ý nghĩa: kernel nhiệt \textbf{không} làm lệch tần số, chỉ \textbf{tắt} biên độ mỗi mode tần số $\beta$ theo hệ số $e^{-\beta^2 t}$. Riêng hằng số ($\beta=0$): tích phân $\tfrac{1}{\sqrt{4\pi t}}\!\int e^{-(x-X)^2/(4t)}dX=1$ (giữ nguyên).

\paragraph{Bước 4 --- tính $A(x,t)$ (mode hằng + mode $\cos2\pi$).}\leavevmode

Hạ bậc: $\sin^2(\pi X)=\tfrac12-\tfrac12\cos(2\pi X)$. Áp bổ đề Bước 3 vào từng mode (tuyến tính):
\[
  A(x,t)=\tfrac12\cdot 1\;-\;\tfrac12\,e^{-(2\pi)^2 t}\cos(2\pi x)
  =\tfrac12\Big(1-e^{-4\pi^2 t}\cos(2\pi x)\Big).
\]

\paragraph{Bước 5 --- tính $B(y,t)$ (mode $\sin$ với $\beta=1$).}\leavevmode

Áp bổ đề với $\beta=1$:
\[
  B(y,t)=e^{-t}\sin(y).
\]

\paragraph{Bước 6 --- ráp lại và kiểm tra.}\leavevmode

Nghiệm thác triển $\bar u=A\cdot B$; hạn chế xuống $R$ (đúng đối xứng) cho \emph{chính} $u$:
\begin{nhobox}
\[
  \boxed{\;u(x,y,t)=\tfrac12\Big(1-e^{-4\pi^2 t}\cos(2\pi x)\Big)\,e^{-t}\sin(y).\;}
\]
\end{nhobox}

\emph{Kiểm tra (ba đầu mối):}
\begin{itemize}
\item \textbf{Dữ kiện đầu} $t=0$: $\tfrac12(1-\cos2\pi x)\sin y=\sin^2(\pi x)\sin y=g$. \checkmark
\item \textbf{Biên Neumann} $x=\pm1$: $u_x=\tfrac12\,e^{-4\pi^2 t}\cdot2\pi\sin(2\pi x)\,e^{-t}\sin y$; tại $x=\pm1$ có $\sin(\pm2\pi)=0$. \checkmark
\item \textbf{Biên Dirichlet} $y=0,\pi$: $\sin(0)=\sin(\pi)=0\Rightarrow u=0$. \checkmark
\end{itemize}
(Theo (a), nghiệm là duy nhất, nên đây chính là nghiệm cần tìm.)

\subsection*{Dạng bài lân cận}
\begin{itemize}
\item \textbf{Nhiệt thanh 1D} (tách biến chuỗi hàm riêng, nghiệm dừng) --- xem mục riêng.
\item \textbf{Nhiệt nửa-dải} ($0<x<\pi$, $y>0$): chuỗi sin theo $x$ $\times$ erf/heat-kernel theo $y$ --- xem mục ``nhiệt miền có biên''.
\item Cặp cạnh Neumann cả hai biến $\Rightarrow$ thác triển chẵn cả hai $\Rightarrow$ chuỗi/cosine; chú ý mode hằng.
\end{itemize}
